\title{Group Sequence Policy Optimization}

\begin{abstract}
We present Group Sequence Policy Optimization (GSPO), a reinforcement learning approach designed for large language model training that is robust, efficient, and delivers strong performance. Distinct from earlier methods that rely on token-level importance ratios, GSPO establishes importance ratios at the sequence level, incorporating sequence-based clipping, rewarding, and optimization processes. Through empirical evaluations, we show that GSPO not only surpasses the GRPO algorithm in both training efficiency and final outcomes but also enhances stability in Mixture-of-Experts (MoE) reinforcement learning contexts. Additionally, GSPO offers opportunities to streamline RL system architecture. These advancements have played a pivotal role in the progress seen with the latest Qwen3 models.
\end{abstract}

\section{Introduction}
\label{sec:intro}

Reinforcement learning (RL) has become a central approach for advancing the scalability of language models \citep{o1,dpsk-r1,qwq32b,qwen3}. By leveraging RL at a large scale, language models acquire the proficiency to address complex challenges—including those found in advanced mathematics and coding—by engaging in more extensive and intricate reasoning.

Scaling up RL via increased computational resources, however, fundamentally depends on ensuring stable and resilient training behaviors. Yet, contemporary leading RL techniques, such as GRPO \citep{grpo}, encounter pronounced instability when employed to train extremely large language models. This often leads to severe, sometimes irreparable, model failures described as collapse \citep{qwen3, minimax-m1}. Such training instability significantly restricts ongoing progress in enhancing language model abilities through prolonged RL optimization.

This study reveals that the root cause of instability in GRPO arises from the improper utilization and breakdown of importance sampling weights within its methodological framework. This flaw leads to the introduction of high-variance noise during training, which grows with response length and is exacerbated by the clipping mechanism, eventually causing the model to collapse.

To overcome these significant challenges, we introduce \textbf{Group Sequence Policy Optimization (GSPO)}, an innovative RL approach tailored for training large language models. GSPO distinguishes itself through a rigorously justified formulation of the importance ratio, drawing upon sequence likelihood as established in \citep{click}, thereby adhering closely to the foundational ideas of importance sampling. Moreover, GSPO determines normalized rewards based on the comparative advantages across multiple responses to a single query, ensuring coherence between sequence-level feedback and the optimization objective.

Our experimental results highlight the marked advantages of GSPO compared to GRPO, particularly regarding training stability, computational efficiency, and achieved performance. Notably, GSPO intrinsically addresses the instability encountered during the reinforcement learning training of large Mixture-of-Experts (MoE) models, removing the necessity for elaborate stabilization techniques and indicating the prospect of a streamlined RL infrastructure. These strengths of GSPO played a crucial role in driving notable performance gains observed in the most recent Qwen3 models. We anticipate that GSPO will serve as a dependable and scalable algorithmic basis for ongoing progress in large-scale reinforcement learning with language models.

\section{Preliminaries}

\paragraph{Notation}
Throughout this work, we refer to an autoregressive language model with parameters $\theta$ as a policy, denoted by $\pi_\theta$. The symbol $x$ represents a query, and $\mathcal{D}$ stands for the set of all queries. For a given query $x$ and its associated response $y$, the probability assigned to $y$ by the policy $\pi_\theta$ is expressed as $\pi_\theta (y | x)=\prod_{t=1}^{|y|} \pi_\theta (y_t | x, y_{<t} )$, where $|y|$ indicates the length of $y$ in tokens. Furthermore, a verifier $r$ can assign a reward value $r(x, y) \in [0, 1]$ to any query-response pair $(x, y)$.

\paragraph{Proximal Policy Optimization (PPO)} Using samples generated from the old policy $\pi_{\theta_\text{old}}$, PPO \citep{ppo} constrains the policy update within a proximal region of the old policy through the clipping mechanism. Specifically, PPO employs the following objective for policy optimization (we omit the KL regularization term hereinafter for brevity, as it is not the focus of this paper):

\begin{align}
\mathcal{J}_\text{PPO}(\theta) = \mathbb{E}_{ x \sim \mathcal{D},\, y \sim \pi_{\theta_\text{old}}( \cdot | x) }
\left[ \frac{1}{|y|} \sum_{t=1}^{|y|} 
\min \left( w_{t}(\theta) \widehat{A}_{t},  \, \mathrm{clip} \left( w_{t}(\theta), 1 - {\varepsilon}, 1 + {\varepsilon}\right) \widehat{A}_{t} \right)
\right],
\end{align}

In this context, the importance ratio for token $y_t$ is given by $w_{t}(\theta) = \frac{ \pi_{\theta} (y_{t} | x, y_{<t}) }{ \pi_{\theta_\text{old}} (y_{t} | x,y_{<t})} $. The advantage $\widehat{A}_{t}$ associated with $y_t$ is computed using a separate value model, while $\varepsilon$ specifies the clipping threshold for the importance ratios.

A significant practical obstacle for PPO arises from its substantial dependence on the value model. In most cases, the value model has approximately the same scale as the policy model, which incurs notable computational and memory costs. The success of the method is further tied to the accuracy of value estimation. Constructing a dependable value model is already a demanding task; extending its robustness to handle longer sequences and increasingly sophisticated tasks amplifies these difficulties.

\paragraph{Group Relative Policy Optimization (GRPO)} GRPO \citep{grpo} bypasses the need for the value model by computing the relative advantage of each response within a group of responses to the same query. Specifically, GRPO optimizes the following objective:

\begin{align}
\label{equ:grpo}
\mathcal{J}_\text{GRPO}(\theta) = \mathbb{E}_{ x \sim \mathcal{D},\, \{y_i\}_{i=1}^G \sim \pi_{\theta_\text{old}}( \cdot | x) }
\left[ \frac{1}{G} \sum_{i=1}^{G} \frac{1}{|y_i|} \sum_{t=1}^{|y_i|} 
\min \left( w_{i,t}(\theta) \widehat{A}_{i,t},  \, \mathrm{clip} \left( w_{i,t}(\theta), 1 - {\varepsilon}, 1 + {\varepsilon}\right) \widehat{A}_{i,t} \right)
\right],
\end{align}

where $G$ is the number of generated responses to each query $x$ (i.e., the group size), and the importance ratio $w_{i,t}(\theta)$ and advantage $\widehat{A}_{i,t}$ of token $y_{i,t}$ are:

\begin{align}
    w_{i,t}(\theta)=\frac{ \pi_{\theta} (y_{i,t} | x, y_{i,<t}) }{ \pi_{\theta_\text{old}} (y_{i,t} | x,y_{i,<t})},\quad\ 
    \widehat{A}_{i,t} = \widehat{A}_{i} = \frac{r(x, y_i) - \mathrm{mean} \left( \{ r(x, y_i) \}_{i=1}^G \right) }{ \mathrm{std} \left( \{ r(x, y_i) \}_{i=1}^G \right) },
\end{align}

respectively, where all the tokens in $y_i$ share the same advantage as $\widehat{A}_{i}$.

\section{Motivation}
\label{sec:motivation}

As models increase in parameter count, exhibit greater sparsity (such as in Mixture-of-Experts architectures), and generate longer outputs, reinforcement learning demands a correspondingly large rollout batch size to ensure optimal use of computational resources. To enhance the efficiency with which data samples are utilized, it is customary to divide this extensive rollout batch into several smaller mini-batches that are then used for performing gradient updates. However, this batching strategy naturally leads to an off-policy learning context: responses $y$ are drawn from a previous version of the policy, denoted by $\pi_{\theta_\text{old}}$, rather than from the current policy $\pi_\theta$ under optimization. This condition highlights the role of the clipping strategies adopted in PPO and GRPO algorithms, which are designed to exclude samples that deviate excessively from the current policy when computing gradients, thereby mitigating the impact of strongly off-policy data.

While mechanisms like clipping aim to manage this off-policy discrepancy, we identify a more fundamental issue in GRPO: \textit{its objective is ill-posed}. This problem becomes particularly acute when training large models on long-response tasks, leading to catastrophic model collapse. The ill-posed nature of the GRPO objective stems from a misapplication of importance sampling weights. The principle of importance sampling is to estimate the expectation of a function $f$ under a target distribution $\pi_\text{tar}$ by re-weighting samples drawn from a behavior distribution $\pi_\text{beh}$:

\begin{align}
\mathbb{E}_{z \sim \pi_\text{tar}} \left[ f(z) \right] = \mathbb{E}_{z \sim \pi_\text{beh}} \left[ \frac{ \pi_\text{tar} (z) }{ \pi_\text{beh} (z)} \, f(z) \right].
\end{align}

A key aspect here is that successful correction for the discrepancy between distributions depends on averaging over a large number of samples ($N \gg 1$) drawn from the behavior policy $\pi_\text{beh}$. Only with sufficient samples does the importance weighting term, $\frac{ \pi_\text{tar} (z) }{ \pi_\text{beh} (z)}$, effectively reconcile differences between the two distributions.

By comparison, GRPO leverages an importance weight of the form $\frac{ \pi_{\theta} (y_{i,t} | x, y_{i,<t}) }{ \pi_{\theta_\text{old}} (y_{i,t} | x,y_{i,<t})}$ at each time step $t$ for every token. This weight, however, is computed using just one sampled token $y_{i,t}$ from the corresponding next-token policy $\pi_{\theta_\text{old}}( \cdot | x, y_{i, <t})$, rendering it ineffective as a correction for off-policy effects. Rather than mitigating distributional bias, it injects substantial variance into gradient estimates during training; such variance can accumulate over lengthy sequences and becomes more severe due to the impact of the clipping procedure. Our experimental findings confirm that this dynamic can precipitate sudden and often unrecoverable model collapse. Even when reverting to earlier checkpoints, rigorously tuning hyperparameters (such as clipping thresholds), lengthening generated outputs, or altering the RL querying approach, training cannot typically be restored once collapse has taken place.

This finding exposes a deeper shortcoming within GRPO's structure. The ineffectiveness of per-token importance weighting underscores a fundamental rule: \textbf{the scale at which optimization is performed should be consistent with the scale of reward assignment}. As the reward functions at the sequence level, employing off-policy adjustments for individual tokens is inherently flawed. This realization leads us to abandon the token-level framework in favor of formulating importance weighting and optimization methods at the level of entire sequences.

\section{Algorithm}

\subsection{GSPO: Group Sequence Policy Optimization}

Although the token-level importance weight $\frac{ \pi_{\theta} (y_{i,t} | x, y_{i,<t}) }{ \pi_{\theta_\text{old}} (y_{i,t} | x,y_{i,<t})}$ poses challenges within GRPO, we note that, for language generation tasks, the \textit{sequence-level} importance weight $\frac{ \pi_{\theta} (y | x) }{ \pi_{\theta_\text{old}} (y | x)}$ possesses an interpretable theoretical justification: it quantifies the extent to which the generated response $y$, drawn from $\pi_{\theta_\text{old}} (\cdot | x)$, differs from the updated model $\pi_{\theta} (\cdot | x)$. This correspondence resonates well with sequence-level reward formulations, and furthermore, offers an intuitive metric for guiding the clipping procedure.

Based on this straightforward observation, we propose the \textbf{Group Sequence Policy Optimization (GSPO)} algorithm. GSPO employs the following sequence-level optimization objective:

\begin{align}
\mathcal{J}_\text{GSPO} (\theta) =
\mathbb{E}_{ x \sim \mathcal{D},\, \{y_i\}_{i=1}^G \sim \pi_{\theta_\text{old}}( \cdot | x) }
\left[ 
\frac{1}{G} \sum_{i=1}^{G}
\min \left( s_{i}(\theta)  \widehat{A}_{i},  \, \mathrm{clip} \left( s_{i}(\theta), 1 - {\varepsilon}, 1 + {\varepsilon} \right) \widehat{A}_{i} \right) 
\right],
\label{equ:gspo}
\end{align}

where we adopt the group-based advantage estimation:

\begin{align}
\widehat{A}_{i} = \frac{r(x, y_i) - \mathrm{mean} \left( \{ r(x, y_i) \}_{i=1}^G \right) }{ \mathrm{std} \left( \{ r(x, y_i) \}_{i=1}^G \right) },
\end{align}

and define the importance ratio $s_{i}(\theta)$ based on sequence likelihood \citep{click}:

\begin{align}
s_{i}(\theta) = \left( \frac{ \pi_{\theta} (y_i | x) }{ \pi_{\theta_\text{old}} (y_i | x)} \right)^{\frac{1}{|y_i|}}
=
\exp \left( \frac{1}{|y_i|} \sum_{t=1}^{|y_i|} \log \frac{ \pi_{\theta} (y_{i,t} | x, y_{i,<t}) }{ \pi_{\theta_\text{old}} (y_{i,t} | x,y_{i,<t})} \right).
\end{align}

Consequently, GSPO performs clipping on full response sequences rather than on individual tokens, in order to filter out highly ``off-policy'' samples when estimating gradients, thereby aligning the process with both sequence-level reward assignment and optimization goals. Importantly, length normalization is utilized in $s_{i}(\theta)$ to help stabilize variance and maintain $s_{i}(\theta)$ within a consistent numerical interval. Failing to normalize length would allow significant changes in a few token probabilities to disproportionately impact the sequence-level importance ratio, necessitating variable clipping thresholds depending on response length. Additionally, it is worth noting that due to their distinct importance ratio definitions, the clipping thresholds used in GSPO generally diverge in magnitude from those found in earlier approaches like GRPO.

\subsection{Gradient Analysis}
\label{subsec:gradient}

We can derive the gradient of the GSPO objective as follows (clipping is omitted for brevity):

\begin{align}
\nabla_{\theta} \mathcal{J}_\text{GSPO} (\theta)
=&\ 
\nabla_{\theta} \mathbb{E}_{ x \sim \mathcal{D},\, \{y_i\}_{i=1}^G \sim \pi_{\theta_\text{old}}( \cdot | x) }
\left[ 
\frac{1}{G} \sum_{i=1}^{G} s_{i}(\theta) \widehat{A}_{i}
\right] \\
= &\ 
\mathbb{E}_{ x \sim \mathcal{D},\, \{y_i\}_{i=1}^G \sim \pi_{\theta_\text{old}}( \cdot | x) }
\left[ 
\frac{1}{G} \sum_{i=1}^{G}
s_{i}(\theta) \widehat{A}_{i}
\cdot \nabla_{\theta} \log s_{i}(\theta)
\right] \\
=&\ 
\mathbb{E}_{ x \sim \mathcal{D},\, \{y_i\}_{i=1}^G \sim \pi_{\theta_\text{old}}( \cdot | x) }
\left[ 
\frac{1}{G} \sum_{i=1}^{G} 
\left( \frac{ \pi_{\theta} (y_i | x) }{ \pi_{\theta_\text{old}} (y_i | x)} \right)^{\frac{1}{|y_i|}} \widehat{A}_{i} 
\cdot \frac{1}{|y_i|} \sum_{t=1}^{|y_i|} \nabla_{\theta} \log \pi_{\theta} (y_{i,t} | x, y_{i,<t}) 
\right].
\label{equ:gspo_grad}
\end{align}

For comparison, the gradient of the GRPO objective is as follows (note that $\widehat{A}_{i,t} = \widehat{A}_{i}$):

\begin{align}
\nabla_{\theta} \mathcal{J}_\text{GRPO} (\theta) 
=&\ 
\nabla_{\theta} \mathbb{E}_{ x \sim \mathcal{D},\, \{y_i\}_{i=1}^G \sim \pi_{\theta_\text{old}}( \cdot | x) }
\left[ \frac{1}{G} \sum_{i=1}^{G} \frac{1}{|y_i|} \sum_{t=1}^{|y_i|} 
w_{i,t}(\theta) \widehat{A}_{i,t}
\right] \\
=&\
\mathbb{E}_{ x \sim \mathcal{D},\, \{y_i\}_{i=1}^G \sim \pi_{\theta_\text{old}}( \cdot | x) }
\left[ \frac{1}{G} \sum_{i=1}^{G} \widehat{A}_{i} 
\cdot \frac{1}{|y_i|} \sum_{t=1}^{|y_i|} 
\frac{ \pi_{\theta} (y_{i,t} | x, y_{i,<t}) }{ \pi_{\theta_\text{old}} (y_{i,t} | x,y_{i,<t})} 
\nabla_{\theta} \log \pi_{\theta} (y_{i,t} | x, y_{i,<t})  
\right].
\end{align}

The key difference between GSPO and GRPO centers on \textit{the method by which they assign weights to the gradients of token log likelihoods}. GRPO applies a per-token weighting based on the so-called ``importance weight'' $\frac{ \pi_{\theta} (y_{i,t} | x, y_{i,<t}) }{ \pi_{\theta_\text{old}} (y_{i,t} | x,y_{i,<t})}$. Notably, these weights are not uniform and can fall within the intervals $(0, 1+\varepsilon]$ when $\widehat{A}_{i} >0$ or $[1-\varepsilon, +\infty)$ when $\widehat{A}_{i} < 0$, which introduces nontrivial variability. Over the course of training, such fluctuating weights may accumulate and introduce undesirable, unpredictable effects. On the other hand, GSPO avoids this problem by applying uniform weights to every token in a response, thereby circumventing the instability that can arise in GRPO.

\subsection{GSPO-token: A Token-level Objective Variant}

In scenarios like multi-turn RL, we may desire a finer-grained advantage adjustment than the sequence level. To this end, we introduce a token-level objective variant of GSPO, namely \textbf{GSPO-token}, to allow token-wise advantage customization:

\begin{align}
\mathcal{J}_\text{GSPO-token}(\theta)
=
\mathbb{E}_{ x \sim \mathcal{D},\, \{y_i\}_{i=1}^G \sim \pi_{\theta_\text{old}}( \cdot | x) }
\left[ 
\frac{1}{G} \sum_{i=1}^{G} \frac{1}{|y_i|} \sum_{t=1}^{|y_i|} 
\min \left( s_{i,t}(\theta) \widehat{A}_{i,t},  \, \mathrm{clip} \left( s_{i,t}(\theta), 1 - {\varepsilon}, 1 + {\varepsilon} \right) \widehat{A}_{i,t} \right) 
\right],
\label{equ:gspo-token}
\end{align}

where

\begin{align}
s_{i,t}(\theta) 
= \mathrm{sg} \left[ s_{i}(\theta) \right]  \cdot \frac{ \pi_{\theta} (y_{i,t} | x, y_{i,<t}) }{ \mathrm{sg} \left[ \pi_{\theta} (y_{i,t} | x, y_{i,<t}) \right] },
\end{align}

and $\mathrm{sg}[\cdot]$ denotes only taking the numerical value but stopping the gradient, corresponding to the \texttt{detach} operation in PyTorch. The gradient of GSPO-token can be derived as:

\begin{align}
\nabla_{\theta} \mathcal{J}_\text{GSPO-token}(\theta)
=&\ 
\nabla_{\theta} \mathbb{E}_{ x \sim \mathcal{D},\, \{y_i\}_{i=1}^G \sim \pi_{\theta_\text{old}}( \cdot | x) }
\left[ 
\frac{1}{G} \sum_{i=1}^{G}
\frac{1}{|y_i|} \sum_{t=1}^{|y_i|} 
s_{i,t}(\theta) \widehat{A}_{i,t}
\right] \\
=&\ 
\mathbb{E}_{ x \sim \mathcal{D},\, \{y_i\}_{i=1}^G \sim \pi_{\theta_\text{old}}( \cdot | x) }
\left[ 
\frac{1}{G} \sum_{i=1}^{G} s_i (\theta) \cdot \frac{1}{|y_i|} \sum_{t=1}^{|y_i|} 
\widehat{A}_{i,t} \frac{ \nabla_{\theta} \pi_{\theta} (y_{i,t} | x, y_{i,<t}) }{ \pi_{\theta} (y_{i,t} | x, y_{i,<t}) }
\right] \\
=&\ 
\mathbb{E}_{ x \sim \mathcal{D},\, \{y_i\}_{i=1}^G \sim \pi_{\theta_\text{old}}( \cdot | x) }
\left[ 
\frac{1}{G} \sum_{i=1}^{G}  \left( \frac{ \pi_{\theta} (y_i | x) }{ \pi_{\theta_\text{old}} (y_i | x)} \right)^{\frac{1}{|y_i|}}
\cdot \frac{1}{|y_i|} \sum_{t=1}^{|y_i|}
\widehat{A}_{i,t} \nabla_{\theta} \log \pi_{\theta} (y_{i,t} | x, y_{i,<t})  
\right].
\label{equ:gspo-token_grad}
\end{align}

Observe that $\frac{ \pi_{\theta} (y_{i,t} | x, y_{i,<t}) }{ \mathrm{sg} \left[ \pi_{\theta} (y_{i,t} | x, y_{i,<t}) \right] }$ evaluates to 1, which implies that $s_{i,t}(\theta)$ is, in effect, equal to $s_{i}(\theta)$ in terms of numerical value.
A comparison of Equation~\eqref{equ:gspo} with~\eqref{equ:gspo-token}, as well as Equation~\eqref{equ:gspo_grad} with~\eqref{equ:gspo-token_grad}, reveals that GSPO-token and GSPO become equivalent both numerically and theoretically (in their optimization objectives, conditions for clipping, and gradient formulas) when each token's advantage in the sequence $y_i$ is assigned the identical value (specifically, $\widehat{A}_{i,t} = \widehat{A}_{i}$). Nevertheless, GSPO-token allows for more nuanced control, enabling advantages to be tuned for each individual token.

\section{Experiments and Discussion}

\subsection{Empirical Results}

We conduct experiments using a cold-start configuration of the model, which is initialized from Qwen3-30B-A3B-Base and fine-tuned for our tasks. We present both the training reward trajectories and performance trends across several benchmarks: AIME'24 (measured as average Pass@1 across 32 samples), LiveCodeBench (for the 202410-202502 split, measured as average Pass@1 across 8 samples), and CodeForces (assessed via Elo Rating). During reinforcement learning training, each batch of generated rollout data is divided into four mini-batches to apply gradient updates. For the GSPO algorithm, the left and right clipping thresholds in Equation~\eqref{equ:gspo} are set to 3e-4 and 4e-4, respectively. For baseline comparison, we use GRPO, configuring its left and right clipping parameters in Equation~\eqref{equ:grpo} to 0.2 and 0.27, respectively; these parameters are finely adjusted for an equitable comparison. It is important to highlight that GRPO requires implementing the Routing Replay training method to ensure stable convergence in MoE RL scenarios—a topic we elaborate on in \S~\ref{subsec:routing_replay}—whereas \textit{GSPO removes the necessity for this approach}.

Figure~\ref{fig:results} illustrates that training using GSPO maintains consistent stability throughout the process. Notably, \textbf{GSPO enables ongoing enhancement in performance as training compute increases, the query set is frequently refreshed, and the generation length is prolonged}. In addition, GSPO outperforms GRPO in terms of training efficiency, reaching higher training accuracy and improved benchmark results for equivalent training compute and query usage. Lastly, GSPO has been effectively integrated into RL training of the most recent Qwen3 models, providing strong evidence for GSPO's ability to harness the benefits of RL scaling within large language models.

\subsection{Curious Observation on Clipping Fractions}

GSPO differs fundamentally from GRPO in that it applies clipping to entire response sequences instead of to individual tokens. As illustrated in Figure~\ref{fig:clipping}, this results in GSPO clipping a proportion of tokens that is two orders of magnitude higher than that of GRPO, regardless of changes to the clipping range. Despite this substantial increase in the number of clipped tokens and, correspondingly, a reduction in the number of tokens available for model updates or gradient calculations, GSPO consistently demonstrates greater training efficiency relative to GRPO. This seemingly paradoxical outcome—where a considerably larger share of tokens are clipped yet training proceeds more efficiently—suggests that the gradient estimates derived at the token level in GRPO are inherently noisy, thereby limiting their utility for making the best use of sampled data. By instead performing clipping at the sequence level, GSPO offers a more stable and informative learning signal, ultimately leading to improved sample utilization.

\subsection{Benefit of GSPO for MoE Training}
\label{subsec:routing_replay}

\paragraph{Background}
In contrast to standard RL training involving dense models, MoE models, with their inherently sparse activations, bring distinct challenges to training stability. Notably, employing the GRPO algorithm exposes issues where \textit{expert-activation volatility} in MoE models can hinder RL optimization from reaching convergence. Specifically, following one or more applications of gradient updates, the set of experts activated for an identical response often varies considerably. For instance, in the 48-layer Qwen3-30B-A3B-Base model, post each RL gradient update, about 10\% of the experts selected for any given rollout sample using the updated policy $\pi_{\theta}$ are not the same as those chosen by the previous policy $\pi_{\theta_\text{old}}$. This effect intensifies with increased model depth, resulting in substantial fluctuations in the token-level importance ratios $w_{i,t}(\theta) = \frac{ \pi_{\theta} (y_{i,t} | x, y_{i,<t}) }{ \pi_{\theta_\text{old}} (y_{i,t} | x,y_{i,<t})}$, as elaborated in \S~\ref{sec:motivation} and~\ref{subsec:gradient}. Such instability further undermines the reliability of these ratios, ultimately impeding effective RL convergence.

\paragraph{Our Previous Approach} In earlier work, we addressed this problem by implementing the \textbf{Routing Replay} training method. This approach involves storing the experts selected by $\pi_{\theta_\text{old}}$ and, during the calculation of the importance weights $w_{i,t}(\theta) = \frac{ \pi_{\theta} (y_{i,t} | x, y_{i,<t}) }{ \pi_{\theta_\text{old}} (y_{i,t} | x,y_{i,<t})}$, using these cached routing decisions within $\pi_{\theta}$. Consequently, for any given token $y_{i,t}$, both $\pi_{\theta} (y_{i,t} | x, y_{i,<t})$ and $\pi_{\theta_\text{old}} (y_{i,t} | x,y_{i,<t})$ operate over an identical set of active experts, thereby maintaining stability in the importance ratios at the token level and facilitating effective optimization of the same activated sub-network throughout gradient-based training. Figure~\ref{fig:routing_replay} illustrates the crucial role that Routing Replay plays in achieving stable and standard convergence within the GRPO training paradigm for MoE models.

\paragraph{Benefit of GSPO} While Routing Replay facilitates the stable training of MoE models with GRPO by ensuring convergence, it introduces additional burdens in terms of memory and communication and may artificially restrict the model’s true capacity. In comparison, as illustrated in Figure~\ref{fig:results}, GSPO entirely removes reliance on Routing Replay by allowing direct computation of the importance ratios $s_i(\theta)$, leading to normal convergence and steady optimization. The primary advantage lies in the observation that GSPO assesses only the sequence-level likelihood (i.e., $\pi_{\theta} (y_i | x)$), showing insensitivity to the per-token likelihood (i.e., $\pi_{\theta} (y_{i,t} | x, y_{i,<t})$). Because the MoE model’s ability to model language remains consistent, significant variations in sequence likelihood do not occur. Consequently, GSPO directly addresses and mitigates the fluctuations in expert activation characteristic of MoE models, rendering complicated solutions like Routing Replay unnecessary. This results in a streamlined, more robust training procedure and allows utilization of the model’s full capacity without imposed limitations.

\subsection{Benefit of GSPO for RL Infrastructure}

Owing to differences in numerical precision between training systems (such as Megatron) and inference systems (such as SGLang and vLLM), it is common practice to utilize the training engine for recalculating the likelihoods of sampled outputs based on the previous policy $\pi_{\theta_\text{old}}$. In contrast, GSPO optimizes over sequence-level likelihoods instead of token-level ones; conceptually, this higher-level perspective exhibits greater robustness to precision mismatches. As a result, GSPO allows for leveraging likelihood values directly from the inference engine during optimization, eliminating the need for additional likelihood computations by the training engine. This direct approach is particularly advantageous in contexts such as partial rollouts, multi-turn RL, and distributed training-inference architectures.

\section{Conclusion}

Group Sequence Policy Optimization (GSPO) is introduced as a novel reinforcement learning framework tailored for large language model training. Building on the concept of importance sampling, GSPO calculates importance ratios through sequence likelihood, implementing clipping, rewarding, and optimization at the sequence level. Empirical results show that GSPO substantially surpasses GRPO in terms of training stability, efficiency, and overall performance, and is especially effective for large-scale reinforcement learning applications involving MoE models. This advancement forms the basis for significant performance gains observed in recent Qwen3 models. With GSPO serving as a robust, scalable algorithmic foundation, further scaling of reinforcement learning is anticipated, promising major advancements in artificial intelligence capabilities.

\end{document}